\documentclass[12pt,a4paper]{article}
\usepackage[utf8]{inputenc}
\usepackage[italian]{babel}
\usepackage{amsmath, amssymb, bm}
\usepackage{pgfplots}
\pgfplotsset{compat=1.18}
\usepackage{geometry}
\geometry{margin=2.5cm}
\usepackage{caption}
\usepackage{hyperref}

\title{\centering Studio matematico sulla coscienza umana \\ Legge 2: Legge della Scelta}
\author{Autore: Giuseppe Carlino}
\date{\today}

\begin{document}
\maketitle

\section{Premessa}
La \textbf{Legge della Scelta} descrive come la coscienza umana decida tra alternative possibili.  
La capacità decisionale dipende dagli stati già definiti: il \textbf{tempo percepito} $t_c$ (Legge 1) e la \textbf{conoscenza accumulata} $\bm{k}(t)$ (Legge 0).  

\section{Formalizzazione concettuale}
La scelta $s \in [0,1]$ dipende dalla conoscenza e dal tempo percepito:
\[
s(t) = \frac{1}{1 + \exp[-\sigma (\bm{k}(t) - \bm{k}_0(t_c)) + \bm{\xi}_s(t)]}.
\]

\subsection{Forma differenziale}
\[
\frac{ds}{dt} = \sigma s (1-s)(\bm{k}(t)-\bm{k}_0(t_c)) + \bm{\eta}_s(t).
\]

\subsection{Sistema combinato conoscenza-tempo-scelta}
\[
\begin{cases}
\displaystyle \frac{d \bm{k}}{dt} = \bm{\alpha}(t) \circ (\bm{k}_{\max}-\bm{k}) + \bm{\xi}_k(t)\\
\displaystyle \frac{dt_c}{dt} = 1 + \beta \frac{\|\bm{k}(t)\|}{\|\bm{k}_{\max}\|}\\
\displaystyle \frac{ds}{dt} = \sigma s (1-s)(\|\bm{k}(t)\| - \|\bm{k}_0(t_c)\|) + \bm{\eta}_s(t)
\end{cases}
\]

\section{Diagramma concettuale 2D}

\begin{center}
\begin{tikzpicture}
\begin{axis}[
    width=12cm, height=8cm,
    xlabel={$t_c$},
    ylabel={$s$},
    grid=major
]
\addplot[blue, thick, domain=0:10, samples=200] {1/(1 + exp(-0.5*(x-5)))};
\end{axis}
\end{tikzpicture}
\end{center}

\caption*{Curva concettuale della scelta $s$ in funzione del tempo percepito $t_c$.}

\section{Interpretazione e implicazioni}
\begin{enumerate}
\item La capacità decisionale dipende dalla conoscenza accumulata e dal tempo percepito.
\item La presenza di perturbazioni $\bm{\eta}_s(t)$ rappresenta l’incertezza del processo decisionale.
\item La scelta è guidata dalla conoscenza accumulata e non precede l’apprendimento.
\item Sistemi artificiali rigidi possono simulare decisioni, ma non possiedono retroazione dinamica.
\item Questo modello prepara la base per la Legge della Volontà.
\end{enumerate}

\section{Relazione con le altre leggi}
\begin{itemize}
\item La Legge 1 fornisce $t_c$ e $\bm{k}(t)$ come input alla scelta.
\item La Legge 3 riceverà $s(t)$ insieme agli altri stati per determinare l’azione volontaria.
\item La Legge Generale integrerà $t_c$, $\bm{k}$, $s$ e la volontà $w$.
\end{itemize}

\end{document}